\documentclass[12pt,a4paper,twoside]{amsbook}
\usepackage[english,french]{babel}
\usepackage{fancyhdr}
\usepackage{amssymb,amsfonts,amsxtra}
\usepackage{ifthen}
\usepackage[all]{xy}
\usepackage{graphicx}
\usepackage{algorithm}
\usepackage{algorithmic}
\usepackage{moreverb}
\usepackage{supertabular}
\usepackage{fancybox}
%\usepackage[first,light,english]{draftcopy}
\includeonly{
abstract,
acknowl,
chap1,
chap2,
chap3,
chap4,
chap5,
appendix,
biblio,
index}

\makeindex

%\parindent=0cm
%\parskip=0.2truecm

\sloppy

\addtolength{\oddsidemargin}{0.5cm}
\addtolength{\evensidemargin}{-0.5cm}

\numberwithin{section}{chapter}
\numberwithin{equation}{section}

%\renewcommand\proofname{{\itshape Proof}}

\swapnumbers
\newtheoremstyle{plain}	%NAME
	{0.6truecm}		%ABOVESPACE
	{}			%BELOWSPACE
	{\itshape}		%BODYFONT
	{\parindent}		%INDENT
	{\bfseries}		%HEADFONT
	{.}			%HEADPUNCT
	{\labelsep}		%HEADSPACE (\newline = linebreak)
	{}			%CUSTOM-HEAD-SPEC
\theoremstyle{plain}
\newtheorem{thm}{Theorem}[section]
\newtheorem*{thm*}{Theorem}
\newtheorem{lem}[thm]{Lemma}
\newtheorem{prop}[thm]{Proposition}
\newtheorem{cor}[thm]{Corollary}
\newtheorem*{cor*}{Corollary}

\newtheoremstyle{definition}%NAME
	{0.6truecm}		%ABOVESPACE
	{}			%BELOWSPACE
	{\normalfont}		%BODYFONT
	{\parindent}		%INDENT
	{\bfseries}		%HEADFONT
	{.}			%HEADPUNCT
	{\labelsep}		%HEADSPACE
	{}			%CUSTOM-HEAD-SPEC
\theoremstyle{definition}
\newtheorem{defn}[thm]{Definition}
\newtheorem*{defn*}{Definition}
\newtheorem{nota}[thm]{Notation}
\newtheorem{conv}[thm]{Convention}
\newtheorem{conj}[thm]{Conjecture}
\newtheorem{quest}[thm]{Question}
\newtheorem{exmp}[thm]{Example}
\newtheorem*{exmp*}{Example}
\newtheorem{alg}[thm]{Algorithm}

\newtheoremstyle{remark}	%NAME
	{0.6truecm}		%ABOVESPACE
	{}			%BELOWSPACE
	{\normalfont}		%BODYFONT
	{\parindent}		%INDENT
	{\bfseries}		%HEADFONT
	{.}			%HEADPUNCT
	{\labelsep}		%HEADSPACE
	{}			%CUSTOM-HEAD-SPEC
\theoremstyle{remark}
\newtheorem{rem}[thm]{Remark}
\newtheorem*{rem*}{Remark}
\newtheorem*{note*}{Note}

\DeclareMathOperator{\id}{id}
\DeclareMathOperator{\proj}{proj}
\DeclareMathOperator{\incl}{incl}
\DeclareMathOperator{\FGC}{FGC}
\DeclareMathOperator{\Hom}{Hom}
\DeclareMathOperator{\map}{map}
\DeclareMathOperator{\red}{red}
\DeclareMathOperator{\rk}{rank}
\DeclareMathOperator{\subker}{sub-ker}
\DeclareMathOperator{\im}{im}
\DeclareMathOperator{\odd}{^{odd}}
\DeclareMathOperator{\even}{^{even}}
\DeclareMathOperator{\teven}{even}
\DeclareMathOperator{\deven}{_{even}}
\DeclareMathOperator{\Tor}{Tor}
\DeclareMathOperator{\bideg}{bideg}
\DeclareMathOperator{\degst}{deg_{st}}
\DeclareMathOperator{\card}{Card}
\DeclareMathOperator{\dual}{dual}
\DeclareMathOperator{\coker}{coker}
\newcommand{\N}{\mathbb{N}}
\newcommand{\Z}{\mathbb{Z}}
\newcommand{\Q}{\mathbb{Q}}
\newcommand{\R}{\mathbb{R}}
\newcommand{\F}{\mathbb{F}}
\newcommand{\A}{\mathcal{A}}
\newcommand{\clearemptydoublepage}{\newpage{\pagestyle{empty}\cleardoublepage}}
\renewcommand{\geq}{\geqslant}
\renewcommand{\leq}{\leqslant}

\begin{document}
\selectlanguage{french}
\frontmatter

\clearemptydoublepage

\begin{titlepage}
{\parindent=0cm\includegraphics[scale = 0.75]{./LogoUNIL.pdf}\hfill\includegraphics[scale = 0.3]{./logoSciences.pdf}}
\vskip0.5truecm
\hrule
\begin{center}
\vskip0.2truecm
{\sc I\ \ n\ \ s\ \ t\ \ i\ \ t\ \ u\ \ t\ \ \ \ d\ \ e\ \ \ \ M\ \ a\ \ t\ \ h\ \ \'e\ \ m\ \ a\ \ t\ \ i\ \ q\ \ u\ \ e\ \ s}
\vskip2truecm
\selectlanguage{english}
{\Large\bf Integral Cohomology of\vskip0.3truecm Finite Postnikov Towers}\\
\selectlanguage{french}
\vskip2truecm
{\Large\bf Th\`ese de doctorat}\\
\vskip0.5truecm
pr\'esent\'ee \`a la\\ 
\ \\
Facult\'e des Sciences de\\
l'Universit\'e de Lausanne\\
\ \\
par\\
\vskip0.5truecm
{\bf Alain Cl\'ement}\\
\vskip0.5truecm
Dipl\^om\'e en Math\'ematiques\\
Universit\'e de Lausanne\\
\vskip2truecm
{\bf Jury}
\vskip0.5truecm
M. le Prof. Bernard Testa, pr\'esident\\
M. le Prof. Dominique Arlettaz, directeur de th\`ese\\
M. le Prof. Carles Casacuberta, expert\\
M. le Prof. J\'er\^ome Scherer, expert\\
\vskip1truecm
LAUSANNE\\
2002
\end{center}
\end{titlepage}
%\thanks{The author was supported in part by FNRS Grant \#000000.}
%\date{\today}
%\maketitle
\selectlanguage{english}

\pagestyle{empty}

%\clearemptydoublepage
%\newpage
%imprimatur
%{\vspace*{-3cm}\hspace*{-5cm}\includegraphics[scale = 0.055]{./Imprimatur.pdf}\vspace*{-5cm}}
%light
\ 
\newpage
%\vskip1.5truecm
\vskip0.5truecm
\begin{flushright}
\emph{
%To My Mom And\\
%To My Grandfather
\ \\
\ }
\end{flushright}
\vskip3truecm
{\small
\begin{quotation}
\ 

\selectlanguage{french}
La derni\`ere d\'emarche de la raison est de reconna\^itre\\ 
qu'il y a une infinit\'e de choses qui la surpassent.
\vskip0.8truecm
\raggedleft Pens\'ees, 267 -- \emph{Blaise Pascal}\\
(1623-1662)

%[\dots] If you can make one heap of all your winnings\\
%And risk it on one turn of pitch-and-toss,\\
%And lose, and start again at your beginnings\\
%And never breathe a word about your loss;\\

%If you can force your heart and nerve and sinew\\
%To serve your turn long after they are gone,\\
%And so hold on when there is nothing in you\\
%Except the Will which says to them: "Hold on"; [\dots]\\

%\vskip0.5truecm
%$$\widetilde{\ \ \ \ \ \ \ \ \ \ \ \ \ \ \ }$$
%\vskip0.5truecm
%\selectlanguage{french}
%[\dots] Si tu peux rencontrer Triomphe apr\`es D\'efaite\\
%Et recevoir ces deux menteurs d'un m\^eme front,\\
%Si tu peux conserver ton courage et ta t\^ete\\
%Quand tous les autres les perdront,\\

%Alors les Rois, les Dieux, la Chance et la Victoire\\
%Seront \`a tout jamais tes esclaves soumis\\
%Et, ce qui vaut bien mieux que les Rois et la Gloire,\\
%Tu seras un Homme, mon fils.\\
\selectlanguage{english}

%\vskip0.8truecm
%\raggedleft If $\sim$ Si -- \emph{Sir Rudyard Kipling}
\end{quotation}}

%\clearemptydoublepage
\newpage
\ 
\newpage

\pagestyle{fancyplain}
\renewcommand{\chaptermark}[1]%
        {\markboth{\sc #1} {}}
\renewcommand{\sectionmark}[1]%
        {\markright{#1}}
\lhead[\fancyplain{}{\thepage}]%
        {\fancyplain{}{\leftmark}}
\rhead[\fancyplain{}{\leftmark}]%
        {\fancyplain{}{\thepage}}
\cfoot[]{}
%\footrulewidth=0,4pt

\setcounter{page}{1}
\tableofcontents

\chapter*{Abstract}

By the work of H. Cartan, it is well known that one can find elements of arbitrarilly high torsion in the integral (co)homology groups of an Eilenberg-MacLane space $K(G,n)$, where $G$ is a non-trivial abelian group and $n\geq 2$.

\medskip
The main goal of this work is to extend this result to H-spaces having more than one non-trivial homotopy group.

\medskip
In order to have an accurate hold on H. Cartan's result, we start by studying the duality between homology and cohomology of $2$-local Eilenberg-MacLane spaces of finite type. This leads us to some improvements of H. Cartan's methods in this particular case.

\medskip
Our main result can be stated as follows. Let $X$ be an H-space with two non-trivial finite $2$-torsion homotopy groups. Then $X$ does not admit any exponent for its reduced integral graded (co)homology group.

\medskip
We construct a wide class of examples for which this result is a simple consequence of a topological feature, namely the existence of a weak retract $X\to K(G,n)$ for some abelian group $G$ and $n\geq2$.

\medskip
We also generalize our main result to more complicated stable \mbox{$2$-stage} Postnikov systems, using the Eilenberg-Moore spectral sequence and analytic methods involving Betti numbers and their asymptotic behaviour.

\medskip
Finally, we investigate some guesses on the non-existence of homology exponents for finite Postnikov towers. We conjecture that Postnikov pieces do not admit any (co)homology exponent.

\medskip
This work also includes the presentation of the ``Eilenberg-MacLane machine'', a C++ program designed to compute explicitely all integral homology groups of Eilenberg-MacLane spaces.

\selectlanguage{french}

\chapter*{R\'esum\'e}

Depuis le s\'eminaire H. Cartan de 1954-55, il est bien connu que l'on peut trouver des \'el\'ements de torsion arbitrairement grande dans l'homologie enti\`ere des espaces d'Eilenberg-MacLane $K(G,n)$ o\`u $G$ est un groupe ab\'elien non trivial et $n\geq2$. 

\medskip
L'objectif majeur de ce travail est d'\'etendre ce r\'esultat \`a des H-espaces poss\'edant plus d'un groupe d'homotopie non trivial.

\medskip
Dans le but de contr\^oler pr\'ecis\'ement le r\'esultat de H. Cartan, on commence par \'etudier la dualit\'e entre l'homologie et la cohomologie des espaces d'Eilenberg-MacLane $2$-locaux de type fini. On parvient ainsi \`a raffiner quelques r\'esultats qui d\'ecoulent des calculs de H. Cartan. 

\medskip
Le r\'esultat principal de ce travail peut \^etre formul\'e comme suit. Soit $X$ un H-espace ne poss\'edant que deux groupes d'homotopie non triviaux, tous deux finis et de $2$-torsion. Alors $X$ n'admet pas d'exposant pour son groupe gradu\'e d'homologie enti\`ere r\'eduite.

\medskip
On construit une large classe d'espaces pour laquelle ce r\'esultat n'est qu'une cons\'equence d'une caract\'eristique topologique, \`a savoir l'existence d'un r\'etract faible $X\to K(G,n)$ pour un certain groupe ab\'elien $G$ et $n\geq2$.

\medskip
On g\'en\'eralise \'egalement notre r\'esultat principal \`a des espaces plus compliqu\'es en utilisant la suite spectrale d'Eilenberg-Moore ainsi que des m\'ethodes analytiques faisant appara\^itre les nombres de Betti et leur comportement asymptotique.

\medskip
Finalement, on conjecture que les espaces qui ne poss\`edent qu'un nombre fini de groupes d'homotopie non triviaux n'admettent pas d'exposant homologique.

\medskip
Ce travail contient par ailleurs la pr\'esentation de la``machine d'Eilenberg-MacLane", un programme C++ con\c{c}u pour calculer explicitement les groupes d'homologie enti\`ere des espaces d'Eilenberg-MacLane.

\selectlanguage{english}

\endinput
\selectlanguage{french}

\chapter*{Remerciements}

Je remercie cordialement \hbox{M. le Prof. Dominique Arlettaz} d'avoir accept\'e de diriger ce travail de th\`ese. Son aide constante et inconditionnelle, ses connaissances, ses qualit\'es humaines et son amiti\'e m'ont soutenu durant toute les phases de ma recherche. Je le remercie de m'avoir donn\'e la possibilit\'e de voyager, de rencontrer de nombreux math\'ematiciens et d'avoir ainsi r\'euni les meilleures conditions pour l'\'elaboration de ce travail. Mes remerciements vont \'egalement \`a \hbox{Mme Chantal Arlettaz}, ainsi qu'\`a \hbox{Mathieu} et \hbox{Vincent}, pour leur hospitalit\'e et leur gentillesse. 

\medskip
Je remercie \hbox{M. le Prof. Carles Casacuberta} ainsi que \hbox{M. le Prof. J\'er\^ome Scherer} d'avoir si spontan\'ement accept\'e d'expertiser ce travail. Il n'aurait jamais abouti sans leurs nombreux conseils, corrections et commentaires qui furent tant d'encouragements. Je tiens finalement \`a remercier \hbox{M. le Prof. Frederick R. Cohen} ainsi que \hbox{M. le Prof. Wojciech Chach\'olski} pour toutes les discussions que nous avons eues et pour leurs excellents conseils.

\medskip
Je remercie chaleureusement \hbox{M. le Prof. Bernard Testa} d'avoir pr\'esid\'e le Jury.

\medskip
Ce travail a \'et\'e partiellement financ\'e par le Fonds National pour la Recherche Scientifique (FNRS).

\selectlanguage{english}

%\aufm{Alain Cl\'ement}

\mainmatter

\renewcommand{\chaptermark}[1]%
        {\markboth{#1} {}}
\renewcommand{\sectionmark}[1]%
        {\markright{\thesection\ #1}}
\lhead[\fancyplain{}{\thepage}]%
        {\fancyplain{}{\rightmark}}
\rhead[\fancyplain{}{\rightmark}]%
        {\fancyplain{}{\thepage}}
\cfoot[]{}
%\footrulewidth=0,4pt

\chapter{Introduction}

Let $X$ be a connected space. One can both consider its graded homotopy group, $\pi_*(X)$, and its graded reduced integral homology group, $\widetilde{H}_*(X;\Z)$. If there exists an integer $h\geq1$ such that $h\cdot\pi_*(X)=0$, we then say that $X$ has a {\it homotopy exponent}. Analogously, if there exists an integer $e\geq1$ such that $e\cdot\widetilde{H}_*(X;\Z)=0$, we then say that $X$ has a {\it homology exponent}. 

The general problem posed by D. Arlettaz is to know whether there is a relationship between homotopy exponents and homology exponents. For instance, is it true that a space with a homotopy exponent has a homology exponent too? In this case, how are these two exponents related, if they are? Or conversely, is it possible for a space without a homotopy exponent to admit a homology exponent?

In the present work, we focus on spaces with homotopy exponents and ask if they have a homology exponent.

The first example of such a space which come to mind is an Eilenberg-MacLane space $K(G,n)$, with $G$ a finite group, abelian if $n\geq2$. 

It is very well known that if $G$ is a finite group, then $K(G,1)$ has a homology exponent. Actually we can prove, by using a transfer argument, that $\card(G)\cdot\widetilde{H}^*(K(G,1);\Z)=0$.

By H. Cartan's work \cite{Ca55}, it is well known that $K(G,n)$ has no homology exponent if $n\geq2$. This result gives a drastic answer to D. Arlettaz's question: the existence of a homotopy exponent does not imply the existence of a homology exponent in general. But is it always the case for $1$-connected spaces?

For instance, and as an obvious consequence of the K\"unneth formula, products of Eilenberg-MacLane spaces, which are called {\it generalized Eilenberg-MacLane spaces} (GEM), neither have a homology exponent.

The purpose of this work is then to investigate the integral (co)homology of spaces which have a finite Postnikov tower, i.e. spaces with finitely many non-trivial homotopy groups (these spaces are also called {\it Postnikov pieces}), and which are not GEM's.

Our main results (see Section \ref{section:main results}) deal with spaces with two non-trivial homotopy groups and with stable $2$-stage Postnikov systems (see p. \pageref{d:two-stage Postnikov system} for a definition). These spaces turn out not to have a homology exponent, although they have a homotopy exponent.

In order to work in an affordable and manageable technical framework, we will only compute at the prime $2$. For simplicity, our results then focus on $2$-local Postnikov pieces of finite type, and most of the time with an H-space structure.
\newpage

\input chap1sect1.tex
\newpage
\input chap1sect2.tex
\newpage
\input chap1sect3.tex
\newpage
\input chap1sect4.tex
\newpage
\input chap1sect5.tex
\chapter{(Co)Homology of Eilenberg-MacLane spaces}\label{chapter:2}

\input chap2sect1.tex
\newpage
\input chap2sect2.tex
\newpage
\input chap2sect3.tex
\newpage
\input chap2sect4.tex
\newpage
\input chap2sect5.tex
\chapter{The ``Eilenberg-MacLane machine''}\label{chapter:EMM}

The ``Eilenberg-MacLane machine'' is a program which computes the homology groups $H_*(K(\Z/2^s,n);\Z)$ for all $n\geq1$, $s\geq1$. 

A first version, written in C, is actually able to do more since it computes $H_*(K(\Z/p^s,n);\Z)$ for all prime $p$. This version is far from being optimized and is based on the ``rude force'' of the computer you are working with. The aim was only to obtain a very reliable code. The program finds, in a very naive way, all admissible $p$-words (see \cite[Expos\'e 9, p. 1]{Ca55} for a definition when $p\not=2$) of given height and degree. It happens that it is a tricky, recursive problem which takes a lot of time and memory to be achieved.

A second version, written in C++, provides an optimized and reliable code at the prime $2$. It is based on the results of Appendix \ref{appendix:elementary results} which mainly state that admissible words are in one-to-one correspondance with admissible sequences. As it is simpler to determine all admissible sequences of given excess and degree than admissible words, the process reveals fast and parcimonious in terms of memory.

Since our interest lies only at the prime $2$ in this work, we will only describe the second version of the ``Eilenberg-MacLane machine'' in Section \ref{section:alg}.

Appendix \ref{a:tables} contains tables which were computed with the ``Eilenberg-MacLane machine''. A glance at these tables reveals some heuristic results that are proved in Section \ref{section:heuristic}.

For those who are interested in the C++ implementation, Appendix \ref{a:C++ implementation} contains some relevant header files.

\input chap3sect1
\newpage
\input chap3sect2
\chapter{Examples of spaces and classification}

\input chap4sect1
\newpage
\input chap4sect2
\newpage
\input chap4sect3
\addtocontents{toc}{\protect\newpage}
\chapter{Cohomology of some $2$-stage Postnikov systems}

\input chap5sect1.tex
\newpage
\input chap5sect2.tex
\newpage
\input chap5sect3.tex
\newpage
\input chap5sect4.tex
\newpage
\input chap5sect5.tex
\newpage
\input chap5sect6.tex
\newpage
\input chap5sect7.tex

\backmatter
{
\appendix

\numberwithin{equation}{chapter}

%\swapnumbers
\theoremstyle{plain}
\newtheorem{Athm}{Theorem}[chapter]
\newtheorem*{Athm*}{Theorem}
\newtheorem{Alem}[Athm]{Lemma}
\newtheorem{Aprop}[Athm]{Proposition}
\newtheorem{Acor}[Athm]{Corollary}
\newtheorem*{Acor*}{Corollary}

\theoremstyle{definition}
\newtheorem{Adefn}{Definition}[chapter]
\newtheorem{Anota}{Notation}[chapter]
\newtheorem{Aconv}{Convention}[chapter]
\newtheorem{Aconj}{Conjecture}[chapter]
\newtheorem{Aexmp}{Example}[chapter]

\theoremstyle{remark}
\newtheorem*{Arem}{Remark}
\newtheorem*{Anote}{Note}
%%%%%%%%%%%%%%%%%%%%%%%%%%%%%%%%
\chapter{Adem relations among Steenrod Squares}\label{a:Adem relations}

\begin{align*}
Sq^iSq^j=\sum_{k=0}^{[i/2]}\binom{j-k-1}{i-2j}Sq^{i+j-k}Sq^k\\
\text{for all $0<i<2j$.}
\end{align*}

\bigskip
\begin{align*}
Sq^1Sq^1=0, Sq^1Sq^3=0,\dots; &&&Sq^1Sq^{2n+1}=0\\
Sq^1Sq^2=Sq^3,Sq^1Sq^4=Sq^5,\dots; &&&Sq^1Sq^{2n}=Sq^{2n+1}\\
Sq^2Sq^2=Sq^3Sq^1, Sq^2Sq^6=Sq^7Sq^1,\dots; &&&Sq^2Sq^{4n-2}=Sq^{4n-1}Sq^1\\
Sq^2Sq^3=Sq^5+Sq^4Sq^1,\dots; &&&Sq^2Sq^{4n-1}=Sq^{4n+1}+Sq^{4n}Sq^1\\
Sq^2Sq^4=Sq^6+Sq^5Sq^1,\dots; &&&Sq^{2}Sq^{4n}=Sq^{4n+2}+Sq^{4n+1}Sq^1\\
Sq^2Sq^5=Sq^6Sq^1,\dots; &&&Sq^2Sq^{4n+1}=Sq^{4n+2}Sq^1\\
Sq^3Sq^2=0, Sq^3Sq^6=0,\dots; &&&Sq^3Sq^{4n+2}=0\\
Sq^3Sq^3=Sq^5Sq^1,\dots; &&&\dots\\
\dots;&&&Sq^{2n-1}Sq^n=0\\
\vdots
\end{align*}
%%%%%%%%%%%%%%%%%%%%%%%%%%%%%%%%%
\input appendix1.tex
%%%%%%%%%%%%%%%%%%%%%%%%%%%%%%%%%
\chapter{Tables for Eilenberg-MacLane spaces homology}\label{a:tables}

The {\bf genus}\index{genus} of a generator $(x,y)$ (see Definition \ref{d:generators}, p. \pageref{d:generators}) is given by
$$
\begin{cases}
1 &\text{if $x$ is of the form $\sigma^n$,}\\
2 &\text{if $x$ is not of the form $\sigma^n$ and $x$ (or $y$) ends with $\sigma$,}\\
3 &\text{if $x$ (or $y$) ends with $\psi_2$.}\\
\end{cases}
$$

\input table1.tex
\newpage
\input table2.tex
\newpage
\input table3.tex
\newpage
\input table4.tex
%%%%%%%%%%%%%%%%%%%%%%%%%%%%%%%%%%%%
\chapter{C++ implementation of the\\ ``Eilenberg-MacLane machine''}\label{a:C++ implementation}

\renewcommand\verbatimtabsize{4\relax}

\section{Header file \texttt{Mod2GradedGroups.hpp}}
{\small
\listinginput[1]{1}{"../Eilenberg-MacLane Machine/Mod2GradedGroups.hpp"}
}

\newpage
\section{Header file \texttt{Mod2AdmissibleSequence.hpp}}
{\small
\listinginput[1]{1}{"../Eilenberg-MacLane Machine/Mod2AdmissibleSequences.hpp"}
}

\newpage
\section{Header file \texttt{Mod2Words.hpp}}
{\small
\listinginput[1]{1}{"../Eilenberg-MacLane Machine/Mod2Words.hpp"}
}

\newpage
\section{Header file \texttt{Mod2ElementaryComplexHomology.hpp}}
{\small
\listinginput[1]{1}{"../Eilenberg-MacLane Machine/Mod2ElementaryComplexHomology.hpp"}
}

\newpage
\section{Header file \texttt{Mod2EilenbergMacLaneHomology.hpp}}
{\small
\listinginput[1]{1}{"../Eilenberg-MacLane Machine/Mod2EilenbergMacLaneHomology.hpp"}
}

\newpage
\section{Header file \texttt{LaTeXLayout.hpp}}
{\small
\listinginput[1]{1}{"../Eilenberg-MacLane Machine/LaTeXLayout.hpp"}
}

}


\endinput
%\chapter*{Notations}

\input mrabbrev
\bibliographystyle{plain}
\begin{thebibliography}{99}

\bibitem{Ad60}{\sc J. F. Adams}, \textit{On the non-existence of elements of Hopf invariant one}, \ANNMAone\ 72 (1960), 20-103.

\bibitem{AAMY01}{\sc D. Arlettaz, C. Ausoni, M. Mimura and N. Yagita}, \textit{The mod $2$ Bockstein spectral sequence for the special linear group over the integers}, \MATPC\ 131 (2001), 445-457.

\bibitem{Bo53}{\sc A. Borel}, \textit{Sur la cohomologie des espaces fibr\'es principaux et des espaces homog\`enes de groupes de Lie compacts}, \ANNMAone\ 57 (1953), 115-207.

\bibitem{Bo54}{\sc A. Borel}, \textit{Sur l'homologie et la cohomologie des groupes de Lie compacts connexes}, \AMEJM\ 76 (1954), 273-342.

\bibitem{Bo55}{\sc A. Borel}, \textit{Topology of Lie groups and characteristic classes}, \BULAMone\ 61 (1955), 397-432.

\bibitem{Bo67}{\sc A. Borel}, \textit{Topics in the homology theory of fibre bundles}, Lecture Notes in Math. 36 (1967).

\bibitem{Bot54}{\sc R. Bott}, \textit{On torsion in Lie groups}, \PRONA\ 40 (1954), 586-588.

\bibitem{Bot56}{\sc R. Bott}, \textit{An application of Morse theory to the topology of Lie groups}, \BULAMone\ 84 (1956), 251-281.

\bibitem{Bot58}{\sc R. Bott}, \textit{The space of loops on a Lie group}, \MICMJ\ 5 (1958), 36-61.

\bibitem{BS53}{\sc R. Bott and H. Samelson}, \textit{On the Pontrjagin product in spaces of paths}, \COMMH\ 27 (1953), 320-337.

\bibitem{BC99}{\sc C. Broto and J. A. Crespo}, \textit{H-spaces with noetherian mod two cohomology algebra}, \TOPOL\ 38 (1999), 353-386.

\bibitem{BCS01}{\sc C. Broto, J. A. Crespo and L. Saumell}, \textit{Non-simply connected H-spaces with finiteness conditions}, \MATPC\ 130 (2001), 475-488.

\bibitem{Br61}{\sc W. Browder}, \textit{Torsion in H-spaces}, \ANNMAone\ 74 (1961), 24-51.

\bibitem{Br63}{\sc W. Browder}, \textit{Higher torsion in H-spaces}, \TRAAM\ 107 (1963), 353-375.

\bibitem{Ca55}{\sc H. Cartan}, \textit{Alg\`ebres d'Eilenberg-MacLane et homotopie}, Expos\'es 2 \`a 16, S\'eminaire Henri Cartan, Ecole Normale Sup\'erieure, Paris, 1956. 

\bibitem{Cl63}{\sc A. Clark}, \textit{On $\pi_3$ of finite dimensional H-spaces}, \ANNMAone\ 78 (1963), 193-196.

\bibitem{Cl00}{\sc A. Cl\'ement}, \textit{Un espace pour lequel $K(\Z/2,2)$ est un r\'etract}, Journal de l'IMA 2 (2000), 3-10. (See {\tt http://www.unil.ch/ima/journal})

\bibitem{Cl01}{\sc A. Cl\'ement}, \textit{Un espace ``sans r\'etract''}, Journal de l'IMA 3 (2001), 39-41. (See {\tt http://www.unil.ch/ima/journal})

\bibitem{CP00}{\sc F. R. Cohen and F. P. Peterson}, \textit{Some free modules over the cohomology of $K(Z/2Z,n)$: a short walk in the Alps}, in: Contemporary Mathematics 265 (2000), 7-20.

\bibitem{DL59}{\sc A. Dold and R. Lashof}, \textit{Principal quasifibrations and fiber homotopy equivalence of bundles}, \ILLJM\ 3 (1959), 285-305.

\bibitem{Dw80}{\sc W. G. Dwyer}, \textit{Higher divided squares in second-quadrant spectral sequences}, \TRAAM\ 260 (1980), 437-447.

\bibitem{DW90}{\sc W. G. Dwyer and C. W. Wilkerson}, \textit{Spaces of null homotopic maps}, Ast\'erisque 191 (1990), 97-108.

\bibitem{EM66}{\sc S. Eilenberg and J. C. Moore}, \textit{Homology and fibrations. I. Coalgebras, cotensor product and its derived functors},  \COMMH\ 40 (1966), 199-236.

\bibitem{FHLT89}{\sc Y. F\'elix, S. Halperin, J.-M. Lemaire and J.-C. Thomas}, \textit{Mod $p$ loop space homology}, \INVEM\ 95 (1989), 247-262.

\bibitem{Gr97}{\sc J. Grodal}, \textit{The transcendence degree of the mod $p$ cohomology of finite Postnikov systems}, preprint (1997).

\bibitem{Ha70}{\sc J. R. Harper}, \textit{On the cohomology of stable two stage Postnikov systems}, \TRAAM\ 152 (1970), 375-388.

\bibitem{Ho40}{\sc H. Hopf}, \textit{\"Uber die Topologie der Gruppen-Mannigfaltigkeiten und Ihre Verallgemeinerungen}, \ANNMAone\ 13 (1940), 22-52.

\bibitem{Ka88}{\sc R. M. Kane}, \textit{The homology of Hopf spaces}, North-Holland mathematical library 40 (Elsevier Science Publishers B.V., 1988).

\bibitem{Kr62}{\sc L. Kristensen}, \textit{On the cohomology of two-stage Postnikov systems}, \ACTAMone\ 107 (1962), 73-123.

\bibitem{Kr63}{\sc L. Kristensen}, \textit{On the cohomology of spaces with two non-vanishing homotopy groups}, \MATHSone\ 12 (1963), 83-105.

\bibitem{KP72}{\sc L. Kristensen and E. K. Pedersen}, \textit{The $\mathcal{A}$-module structure for the cohomology of two-stage spaces}, \MATHSone\ 30 (1972), 95-106.

\bibitem{LS86}{\sc J. Lannes and L. Schwartz}, \textit{\`A propos de conjectures de Serre et Sullivan}, \INVEM\ 83 (1986), 593-603.

\bibitem{LS89}{\sc J. Lannes and L. Schwartz}, \textit{Sur les groupes d'homotopie des espaces dont la cohomologie modulo $2$ est nilpotente}, \ISRJM\ 66 (1989), 260-273.

\bibitem{Le95}{\sc R. Levi}, \textit{On finite groups and homotopy theory}, Memoirs of the A.M.S. 567, vol. 118, 1995.

\bibitem{Li87}{\sc J. P. Lin}, \textit{Steenrod connections and connectivity in H-spaces}, Memoirs of the A.M.S. 369, vol. 68, 1987.

\bibitem{Mc00}{\sc J. McCleary}, \textit{A user's guide to spectral sequences}, Cambridge studies in advanced mathematics 58, 2nd edition (Cambridge Univ. Press, 2000).

\bibitem{MN84}{\sc C. A. McGibbon and J. A. Neisendorfer}, \textit{On the homotopy groups of a finite dimensional space}, \COMMH\ 59 (1984), 253-257.

\bibitem{Milg66}{\sc R. J. Milgram}, \textit{The bar construction and abelian H-spaces}, \ILLJM\ 11 (1967), 242-250.

\bibitem{Milg69}{\sc R. J. Milgram}, \textit{The structure over the Steenrod algebra of some $2$-stage Postnikov systems}, \QUAJMthree\ 20 (1969), 161-169.

\bibitem{Mi84}{\sc H. Miller}, \textit{The Sullivan conjecture on maps from classifying spaces}, \ANNMAone\ 120 (1984), 39-87.

\bibitem{MM65}{\sc J. W. Milnor and J. C. Moore}, \textit{On the structure of Hopf algebras}, \ANNMAone\ 81 (1965), 211-264.

\bibitem{MT91}{\sc M. Mimura and H. Toda}, \textit{Topology of Lie groups I and II}, Translations of Math. Monographs 91 (1991).

\bibitem{MT68}{\sc R. E. Mosher and M. C. Tangora}, \textit{Cohomology operations and applications in homotopy theory}, Harper's Series in Modern Mathematics (Harper \& Row Publishers Inc., 1968).

\bibitem{Mo56}{\sc J. C. Moore}, \textit{Semi-simplicial complexes and Postnikov systems}, Proc. Internat. Sympos. Algebraic Topology (Mexico, 1956), Universidad Nacional Aut\'onoma de M\'exico and UNESCO, Mexico City, 1958.

\bibitem{MS68-I}{\sc J. C. Moore and L. Smith}, \textit{Hopf algebras and multiplicative fibrations, I}, \AMEJM\ 90 (1968), 752-780.

\bibitem{MS68-II}{\sc J. C. Moore and L. Smith}, \textit{Hopf algebras and multiplicative fibrations, II}, \AMEJM\ 90 (1968), 1113-1150.

\bibitem{OY86}{\sc N. Oda and Z.-Y. Yosimura}, \textit{On the McGibbon and Neisendorfer theorem resolving the Serre conjecture}, \MEMFS\ 40 (1986), 125-135.

\bibitem{Po96}{\sc N. Pointet-Tischler}, \textit{Invariants de Postnikov des espaces de lacets}, Ph.D. thesis, University of Lausanne, Switzerland, 1996.

\bibitem{Re70}{\sc D. Rector}, \textit{Steenrod operations in the Eilenberg-Moore spectral sequence}, \COMMH\ 45 (1970), 540-552.

\bibitem{Sc71}{\sc C. Schochet}, \textit{A two stage Postnikov system where $E_2\not=E_\infty$ in the Eilenberg-Moore spectral sequence}, \TRAAM\ 157 (1971), 113-118.

\bibitem{Schw94}{\sc L. Schwartz}, \textit{Unstable modules over the Steenrod algebra and Sullivan's fixed point set conjecture}, Chicago Lectures in Mathematics Series (The University of Chicago Press, 1994).

\bibitem{Se51}{\sc J.-P. Serre}, \textit{Homologie singuli\`ere des espaces fibr\' es}, \ANNMAone\ 54 (1951), 425-505.

\bibitem{Se53}{\sc J.-P. Serre}, \textit{Cohomologie modulo $2$ des complexes d'Eilenberg-MacLane}, \COMMH\ 27 (1953), 198-232.

\bibitem{Si73-I}{\sc W. M. Singer}, \textit{Steenrod squares in spectral sequences, I}, \TRAAM\ 175 (1973), 327-336.

\bibitem{Si73-II}{\sc W. M. Singer}, \textit{Steenrod squares in spectral sequences, II}, \TRAAM\ 175 (1973), 337-353.

\bibitem{Sm67}{\sc L. Smith}, \textit{Cohomology of stable two stage Postnikov systems}, \ILLJM\ 11 (1967), 310-329.

\bibitem{Sm70}{\sc L. Smith}, \textit{Lectures on the Eilenberg-Moore spectral sequence}, Lecture Notes in Mathematics 134 (1970).

\bibitem{SE62}{\sc N. E. Steenrod and D. B. A. Epstein}, \textit{Cohomology operations}, Annals of Math. Studies 50, Princeton Univ. Press, 1962.
1986, pp. 34--56.

\bibitem{St97}{\sc B. Stroustrup}, \textit{The C++ programming language}, 3rd edition (Addison Wesley Publishing Company, 1997).

\bibitem{To62}{\sc E. Thomas}, \textit{On the mod $2$ cohomology of certain H-spaces}, \COMMH\ 37 (1962), 132-140.

\bibitem{We94}{\sc S. Wenhuai}, \textit{Note on the Serre conjecture theorem of McGibbon and Neisendorfer}, \TOPOA\ 55 (1994), 195-202.

\bibitem{Za76}{\sc A. Zabrodsky}, \textit{Hopf spaces}, North-Holland Math. Study 22 (North-Holland Publishing Company, 1976).

\end{thebibliography}
%\chapter*{Index}
\printindex % ajouter \tiny dans le fichier main.ind

%%%%%%%%%%%%%%%%%%%%%%%%%

\end{document}